\Article{4}
It is now easy to see what magnetic force the Earth exerts at each point in space, however the magnetic fluids within it may be distributed. Imagine the entire volume of the Earth, insofar as it contains free magnetism, that is, separated magnetic fluids, divided into infinitely small elements. Let $\mathrm{d}\mu$ denote in general the quantity of free magnetic fluid in each element, with southern fluid always taken as negative. Let $\rho$ be the distance of $\mathrm{d}\mu$ from an arbitrary point in space with rectangular coordinates $x,y,z$. Finally, let $V$ be the sum of the terms $\frac{\mathrm{d}\mu}{\rho}$ with sign reversed, taken over all the magnetic particles of the Earth. or let
\SourceLeaf{132}
\[
V=-\int\frac{\mathrm{d}\mu}{\rho}
\]
Thus $V$ has a definite value at each point in space. It is a function of $x,y,z$, or of any other three variables by which the points of space are distinguished. The magnetic force $\psi$ at each point in space, and the components $\xi,\eta,\zeta$ obtained by resolving $\psi$ parallel to the coordinate axes, are then given by the formulas
\[
\xi=\frac{\mathrm{d}V}{\mathrm{d}x},\quad \eta=\frac{\mathrm{d}V}{\mathrm{d}y},\quad \zeta=\frac{\mathrm{d}V}{\mathrm{d}z},\quad \psi=\BareRoot{(\xi\xi+\eta\eta+\zeta\zeta)}
\]

\Article{5}
The first step is to derive some general propositions independent of the form of the function $V$, remarkable for their simplicity and elegance.

The total differential of $V$ is
\[
\begin{aligned}
\mathrm{d}V&=\frac{\mathrm{d}V}{\mathrm{d}x}\cdot\mathrm{d}x+\frac{\mathrm{d}V}{\mathrm{d}y}\cdot\mathrm{d}y+\frac{\mathrm{d}V}{\mathrm{d}z}\cdot\mathrm{d}z\\
&=\xi\,\mathrm{d}x+\eta\,\mathrm{d}y+\zeta\,\mathrm{d}z
\end{aligned}
\]

Let $\mathrm{d}s$ denote the distance between the two points to which $V$ and $V+\mathrm{d}V$ refer, and $\theta$ the angle that the direction of the magnetic force $\psi$ makes with $\mathrm{d}s$. Then
\[
\mathrm{d}V=\psi\cos\theta\cdot\mathrm{d}s
\]
because $\frac{\xi}{\psi},\frac{\eta}{\psi},\frac{\zeta}{\psi}$ are the cosines of the angles that the direction of $\psi$ makes with the coordinate axes, while $\frac{\mathrm{d}x}{\mathrm{d}s},\frac{\mathrm{d}y}{\mathrm{d}s},\frac{\mathrm{d}z}{\mathrm{d}s}$ are the cosines of the angles between $\mathrm{d}s$ and the same axes. Thus $\frac{\mathrm{d}V}{\mathrm{d}s}$ equals the force projected onto the direction of $\mathrm{d}s$. The same follows directly from the equation $\frac{\mathrm{d}V}{\mathrm{d}x}=\xi$ once we recall that the coordinate axes can be chosen arbitrarily.

\Article{6}
Suppose two points in space, $P^0,P'$, are joined by an arbitrary line, with $\mathrm{d}s$ representing an unspecified element. Let $\theta$, as before, denote the angle between $\mathrm{d}s$ and the direction of the \SourceLeaf{133}magnetic force there, and $\psi$ its intensity. Then
\[
\int\psi\cos\theta\cdot\mathrm{d}s=V'-V^0
\]
if the integration is taken along the entire line and $V^0,V'$ denote the values of $V$ at the endpoints.

The following corollaries of this fruitful proposition deserve special mention.

I. The integral $\int\psi\cos\theta\cdot\mathrm{d}s$ has the same value along any path from $P^0$ to $P'$.

II. The integral $\int\psi\cos\theta\cdot\mathrm{d}s$, taken along the entire length of any line that returns to its starting point, is always $=0$.

III. Along a closed line, unless $\theta=90^0$ throughout, some of the values of $\theta$ must be less and some greater than $90^0$.

\Article{7}
\expandafter\def\csname HistoricalNoteText15\endcsname{If the function $V$ could have any arbitrarily chosen form, then in special cases a maximum or minimum value of $V$ could correspond to an isolated point or an isolated line, with only smaller or only larger values all around it. It could also correspond to a surface with smaller or larger values on \textit{both} sides at once. The conditions imposed on the function $V$, however, rule out these exceptional cases. A full treatment of the subject is unnecessary for our present purpose and must be left for another occasion.}
The surface on which $V$ has the same definite value $=V^0$ at every point separates the points of space at which $V$ is greater than $V^0$ from those at which it is less than $V^0$\HistoricalNote{15}. It follows readily from the proposition in Article 5 that at each point of this surface the magnetic force is perpendicular to the surface, directed toward the side on which the larger values of $V$ occur. Let $\mathrm{d}s$ be an infinitely small line perpendicular to the surface and $V^0+\mathrm{d}V^0$ the value of $V$ at its other endpoint. Then the intensity of the magnetic force is $=\frac{\mathrm{d}V^0}{\mathrm{d}s}$. All the points at which $V=V^0+\mathrm{d}V^0$ form a second surface infinitely close to the first. At the various places throughout the intervening space, the intensity of the magnetic force is inversely proportional to the distance between the two surfaces. If $V$ is made to vary by infinitely small but equal steps, this gives a system \SourceLeaf{134}of surfaces dividing space into infinitely thin layers. The inverse proportionality of layer thickness to magnetic force intensity then holds at different places in one and the same layer as well as across different layers.

\Article{8}
We now consider how the values of $V$ behave on the Earth's surface.

At a point $P$ on the Earth's surface, let $\psi$ be the intensity and $PM$ the direction of the total magnetic force. Let $\omega$ be the intensity and $PN$ the direction of the force projected onto the horizontal plane. Equivalently, $PN$ is the direction of the magnetic meridian, taken from the south pole of the magnetic needle toward its north pole. Let $i$ be the angle between $PM$ and $PN$, or the inclination. Let $\theta,t$ be the angles between the element $\mathrm{d}s$ of a line lying on the Earth's surface and the directions $PM,PN$. Finally, let $V$ and $V+\mathrm{d}V$ correspond to the starting point and endpoint of $\mathrm{d}s$. We therefore have
\[
\cos\theta=\cos i\cos t,\qquad \omega=\psi\cos i
\]
and the equation in Article 5 becomes
\[
\mathrm{d}V=\omega\cos t\cdot\mathrm{d}s
\]

Suppose, then, that two points $P^0,P'$ on the Earth's surface, at which $V$ has the values $V^0,V'$, are joined by a line lying entirely on the Earth's surface, and that $\mathrm{d}s$ denotes an unspecified element of this line. Then
\[
\int\omega\cos t\cdot\mathrm{d}s=V'-V^0
\]
if the integration is taken along the entire line. Clearly, three corollaries quite similar to those given in Article 6 also apply here.

I. The integral $\int\omega\cos t\,\mathrm{d}s$ has the same value along any path on the Earth's surface from $P^0$ to $P'$.

II. The integral $\int\omega\cos t\cdot\mathrm{d}s$ along the entire length of a closed line lying on the Earth's surface is always $=0$.

III. Along such a closed line, unless \SourceLeaf{135}$t=90^0$ throughout, some of the values of $t$ must necessarily be acute and some obtuse.

\Article{9}
Propositions I and II of the preceding article (which are really only two different ways of expressing the same thing) can be tested, at least approximately, against actual observations. Let $P^0P'P''....P^0$ be a polygon on the Earth's surface whose sides are the shortest lines between their endpoints. Thus, if we treat the Earth here as spherical, the sides are arcs of great circles. Let $\omega^0,\omega',\omega''$ and so on be the intensities of the horizontal magnetic force at the points $P^0,P',P''$ and so on. Further, let $\delta^0,\delta',\delta''$ and so on be the declinations, which we may take in the usual way as positive west of north and negative east of north. Finally, let $(01)$ be the azimuth of the line $P^0P'$ at $P^0$, measured in the usual way from south in the direction of west. Likewise, let $(10)$ be the azimuth of the same line taken in reverse at $P'$, and so on.

Note. that $t$ changes continuously along each side of the polygon but jumps at the vertices, so that it has two different values there. For example, at $P'$ the quantity $t$ has

the value $(10)+\delta'$ when $P'$ is the endpoint of $P^0P'$,

the value $180^0+(12)+\delta'$ when $P'$ is the starting point of $P'P''$.

As an approximate value for the integral $\int\omega\cos t\cdot\mathrm{d}s$ along $P^0P'$, we can take
\[
\tfrac12(\omega^0\cos t^0+\omega'\cos t')\cdot P^0P'
\]
where $t^0,t'$ denote the values of $t$ at $P^0$ as the starting point and at $P'$ as the endpoint of $P^0P'$. This approximation is all that can be obtained insofar as the values of $\omega$ and $t$ are available only at the endpoints $P^0,P'$. The shorter the line, the more justified the approximation. In our notation, the expression given is
\[
=\tfrac12\{\omega'\cos((10)+\delta')-\omega^0\cos((01)+\delta^0)\}\cdot P^0P'
\]
Similarly, the approximate value of the integral along $P'P''$ is
\SourceLeaf{136}
\[
=\tfrac12\{\omega''\cos((21)+\delta'')-\omega'\cos((12)+\delta')\}\cdot P'P''
\]
and so on around the entire polygon.

For a triangle, then, our proposition gives the approximately correct equation
\[
\begin{aligned}
&\omega^0\{P^0P'\cos((01)+\delta^0)-P^0P''\cos((02)+\delta^0)\}\\
{}+{}&\omega'\{P'P''\cos((12)+\delta')-P^0P'\cos((10)+\delta')\}\\
{}+{}&\omega''\{P^0P''\cos((20)+\delta'')-P'P''\cos((21)+\delta'')\}\\
={}&0
\end{aligned}
\]

Clearly, the units of intensity and distance in this equation are arbitrary.

